\section{Scope and Objective}

\subsection{What this document is trying to do}

The repository is not building a naive ``trade from one indicator'' bot. The intended end state is a research and trading framework that studies how professional intraday traders use \Important{visible liquidity}, \Important{hidden liquidity}, \Important{tape aggression}, \Important{context}, and \Important{risk discipline} to decide when a pattern is worth trading and when it is not. The materials in the repository are discretionary and educational, not machine-ready. Therefore the first job is formalization: to translate the trading ideas into explicit statements about what is being observed, what is being inferred, what is likely edge versus narrative, and what can later be tested in data.

This document treats the source materials as a layered knowledge base rather than a single coherent system. The books provide broader framing about scalping, prop infrastructure, and risk management. The transcript series is richer in discretionary pattern language: densities, icebergs, robots, springs, breakouts, false moves, news reactions, and trade review. The practical task is to unify those into one research-oriented reference.

\subsection{Primary repository direction encoded here}

The document is written to support the following repository direction:

\begin{itemize}[leftmargin=1.5em]
\item Work from order book / DOM / level 2, tape / prints, and execution feedback together rather than from order book snapshots alone.
\item Detect patterns explicitly first, before attempting machine learning.
\item Treat the old Russian historical order book data as auxiliary and weak, not as a gold-standard training set.
\item Build toward logging, normalization, replay, feature extraction, setup detection, analytics, paper trading, and later live execution.
\item Include \Important{day-specific adaptation}: the system should learn what appears to be working \emph{today} for \emph{this instrument}, instead of assuming setup quality is constant across sessions.
\end{itemize}

\subsection{Interpretive stance}

The source material speaks in the language of experienced traders. That language contains both hard observations and soft judgment. A recurring pattern throughout the material is that the same visible structure can mean opposite things in different contexts. For example, a large density can mean support, passive absorption, a fake wall, an invitation to a breakout, or bait laid by a robot. Therefore the correct formalization is not ``pattern X always means Y,'' but rather:

\begin{enumerate}[leftmargin=1.5em]
\item identify what is objectively visible,
\item identify what the trader is hypothesizing about other participants,
\item identify what contextual evidence upgrades or downgrades that hypothesis,
\item specify the invalidation condition and the trade management consequences.
\end{enumerate}

That is the structure used throughout the rest of this document.
